\section{Lambda Generators for Diagonal Circuits}
%
\label{App:DiagGate}

\newcommand{\One}{\mathsf{One}}

This section shows how to construct the $\lambda( x, \alpha )$ circuits as described in~\cref{Sect:Hierarchical}.
It is also shown that these circuits belong to the sub-language $\Core$ as defined in in~\cref{Sect:Completeness}.
Up to a change of basis, every $\lambda( x, \alpha )$ generator can be thought of as an $n$-fold controlled global phase by an angle of $\alpha\pi$.
The aforementioned change of basis should apply $X$-gates to the qubits, so that the state $\ket{x} := \bigotimes_{j=0}^{n-1} \ket{x_j}$ is sent to the state $\ket{\One( n )} := \bigotimes_{j=1}^n \ket{1}$.

\subsection{Constructing the Change of Basis}

\newcommand{\Nots}{\mathsf{Nots}}

Let $x \in \{ 0, 1 \}^n$.
Then define $\Nots( x ) = \bigboxtimes_{j=0}^{n-1} \left( \scalebox{0.65}{\circimg{tikz-cache/diag__not.pdf}} \right)$.
Note that in $\Nots( x )$, each $X$-gate has been exponentiated by either $0$ or $1$.
This ensures that $\Nots( x ) \in \Core( n, n )$.
It follows that $\sem{\Nots( x )}_H = \bigotimes_{j=0}^{n-1} X^{1-x_j}$.
To see how this operation acts on $\ket{x}$, there are two cases to consider.
\begin{itemize}
\item If $x_j = 0$, then $X^{1-x_j} \ket{x_j} = X^1 \ket{0} = \ket{1}$.
\item If $x_j = 1$, then $X^{1-x_j} \ket{x_j} = X^0 \ket{1} = \ket{1}$.
\end{itemize}
In either case, $X^{1-x_j} \ket{x_j} = \ket{1}$.
This means that $\sem{\Nots( x )}_H \ket{x} = \bigotimes_{j=0}^{n-1} X^{1-x_j} \ket{x_j} = \ket{\One( n )}$, as desired.
Since $\sem{\Nots( x )}_H$ is unitary, then $\sem{\Nots( x )}_H \ket{\psi} = \ket{\One( n )}$ if and only if $\ket{\psi} = \ket{x}$.
Moreover, $\sem{\Nots( x )}_H$ is self-inverse.

\subsection{Constructing the Controlled Phases}

\newcommand{\Phase}{\mathsf{CPh}}

Let $n \in \mathbb{N}$ and $\alpha \in \mathbb{R}$.
Then the $n$-fold controlled phase gate can be defined inductively as follows.
\begin{align*}
    \Phase( \alpha, 0 )
    &=
    \scalebox{0.65}{\circimg{tikz-cache/diag__phase0.pdf}}
    &
    \Phase( \alpha, k + 1 )
    &=
    \scalebox{0.65}{\circimg{tikz-cache/diag__phasek.pdf}}
\end{align*}
Since $\Phase( \alpha, 0 )$ is a generator for $\Core$ which is closed under the application of $Z$-controls, then $\Phase( \alpha, n ) \in \Core( n, n )$ for each $n \in \mathbb{N}$.
It follows by induction that the following equation holds.
\begin{equation}
    \label{Eqn:DiagHyp}
    \sem{\Phase( \alpha, n )}_H \ket{y} = \begin{cases}
        e^{i\alpha} \ket{y} & \text{if } \ket{y} = \ket{\One( n )} \\
        \ket{y} & \text{otherwise}
    \end{cases}
\end{equation}
The proof proceeds as follows.
\begin{itemize}
\item \textbf{Base Case}.
      Since $\ket{\One( 0 )}$ is the only basis vector for $\mathbb{C}^1$ and $\Phase( \alpha, 0 ) = e^{i\alpha}$, then \cref{Eqn:DiagHyp} holds when $n = 0$.
\item \textbf{Inductive Hypothesis}.
      The equation \cref{Eqn:DiagHyp} holds for some $k \in \mathbb{N}$.
\item \textbf{Inductive Step}.
      Assume that the inductive hypothesis holds for some $k \in \mathbb{N}$.
      This means that \cref{Eqn:DiagHyp} holds for $k$.
      To show that \cref{Eqn:DiagHyp} holds for $k + 1$, let $x \in \{ 0, 1 \}^{k+1}$.
      Then there exists some $a \in \{ 0, 1 \}$ and $y \in \{ 0, 1 \}^k$ such that $x = ay$.
      This means that $\ket{x} = \ket{a} \otimes \ket{y}$.
      Now recall from \cref{Thm:GenCtrlNF} that $\sem{\Phase( \alpha, k + 1 )}_H = \ket{0}\bra{0} \otimes \mathbb{I}_{2^k} + \ket{1}\bra{1} \otimes \sem{\Phase( \alpha, k )}_H$.
      There are three cases to consider.
      \begin{enumerate}
      \item Assume that $a = 0$.
            Then $\sem{\Phase( \alpha, k + 1 )}_H \ket{x} = \ket{a} \otimes \ket{y} = \ket{x}$, since $\braket{0|a} = 1$ and $\braket{1|a} = 0$.
            Then $\ket{x} \ne \ket{\One( k + 1 )}$ and $\sem{\Phase( \alpha, k + 1 )} \ket{x} = \ket{x}$.
      \item Assume $a = 1$ and $\ket{y} \ne \ket{\One( k )}$.
            Then $\sem{\Phase( \alpha, k + 1 )}_H \ket{x} = \ket{a} \otimes \sem{\Phase( \alpha, k )} \ket{y}$, since $\braket{0|a} = 1$ and $\braket{1|a} = 0$.
            Since \cref{Eqn:DiagHyp} holds for $k$ and $\ket{y} \ne \ket{\One( k )}$, then $\sem{\Phase( \alpha, k )} \ket{y} = \ket{y}$.
            It follows that $\sem{\Phase( \alpha, k + 1 )} \ket{x} = \ket{a} \otimes \ket{y} = \ket{x}$.
            In conclusion, $\ket{x} \ne \ket{\One( k + 1 )}$ and $\sem{\Phase( \alpha, k + 1 )} \ket{x} = \ket{x}$.
      \item Assume $a = 1$ and $\ket{y} = \ket{\One( k )}$.
            Then $\sem{\Phase( \alpha, k + 1 )}_H \ket{x} = \ket{a} \otimes \sem{\Phase( \alpha, k )} \ket{y}$, since $\braket{0|a} = 1$ and $\braket{1|a} = 0$.
            Since \cref{Eqn:DiagHyp} holds for $k$ and $\ket{y} = \ket{\One( k )}$, then $\sem{\Phase( \alpha, k )} \ket{y} = e^{i\alpha} \ket{y}$.
            It follows that $\sem{\Phase( \alpha, k + 1 )} \ket{x} = \ket{a} \otimes e^{i\alpha} \ket{y} = e^{i\alpha} \ket{x}$.
            In conclusion, $\ket{x} \ne \ket{\One( k + 1 )}$ and $\sem{\Phase( \alpha, k + 1 )} \ket{x} = e^{i\alpha} \ket{x}$.
      \end{enumerate}
      Then \cref{Eqn:DiagHyp} holds for $k + 1$ by case distinction.
\end{itemize}
Then by the principle of induction, \cref{Eqn:DiagHyp} holds for each $n \in \mathbb{N}$.

\subsection{Constructing the Lambda Generators}

Let $x \in \{ 0, 1 \}^*$ and $\alpha \in (-\pi, \pi]$.
Since $x \in \{ 0, 1 \}^*$, then $x$ is a possibly empty word over the alphabet $\{ 0, 1 \}$.
The \emph{$(x, \alpha)$-lambda generator} is the circuit $\lambda( x, \alpha ) = \Nots( x ) \circ \Phase( \alpha, n ) \circ \Nots( x )$ where $n = |x|$.
Since $\Nots( x ) \in \Core( n, n )$ and $\Phase( \alpha, n ) \in \Core( n, n )$, then $\lambda( x, \alpha ) \in \Core( n, n )$ as well.
It must be shown that $\lambda( x, \alpha )$ has the intended semantics.
Clearly $\{ \ket{y} : y \in \{ 0, 1 \}^n \}$ is a basis for $\mathbb{C}^{2^n}$.
Let $y \in \{ 0, 1 \}^n$.
There are two cases to consider.
\begin{enumerate}
\item Assume that $\ket{y} = \ket{x}$.
      Then the following equation holds.
      \begin{align*}
          \sem{\lambda( x, \alpha )}_H \ket{y}
          &=
          \sem{\Nots( x )}_H \circ \sem{\Phase( \alpha, n )}_H \circ \sem{\Nots( x )}_H \ket{y}
          \\
          &=
          \sem{\Nots( x )}_H \circ \sem{\Phase( \alpha, n )}_H \ket{\One( n )}
          \\
          &=
          e^{i\alpha\pi} \sem{\Nots( x )}_H \ket{\One( n )}
          =
          e^{i\alpha} \ket{x}
      \end{align*}
      Then $\ket{y}$ is an eigenvector of $\sem{\lambda( x, \alpha)}_H$ with eigenvalue $e^{i\alpha}$.
\item Assume that $\ket{y} = \ket{x}$.
      Then there exists some $z \in \{ 0, 1 \}^n$ such that $\ket{z} = \sem{\Nots( x )}_H \ket{y}$ with $z \ne x$.
      Then the following equation holds.
      \begin{align*}
          \sem{\lambda( x, \alpha )}_H \ket{y}
          &=
          \sem{\Nots( x )}_H \circ \sem{\Phase( \alpha, n )}_H \circ \sem{\Nots( x )}_H \ket{y}
          \\
          &=
          \sem{\Nots( x )}_H \circ \sem{\Phase( \alpha, n )}_H \ket{z}
          \\
          &=
          \sem{\Nots( x )}_H \ket{z}
          =
          \ket{y}
      \end{align*}
      Then $\ket{y}$ is an eigenvector of $\sem{\lambda( x, \alpha)}_H$ with eigenvalue $1$.
\end{enumerate}
Then $\sem{\lambda( x, \alpha )}_H$ has eigenvalue $e^{i\alpha}$ with eigenspace $\spn\{ \ket{x} \}$ and eigenvalue $1$ with eigenspace $\spn\{ \ket{y} : y \in \{ 0, 1 \}^n \setminus \{ x \} \}$.
Since $\alpha \in (-\pi, \pi]$, then $\Log\sem{\lambda( x, \alpha )}_H = i \alpha \ket{x} \bra{x} = i \alpha \bigotimes_{j=0}^{n-1} \ket{x_j} \bra{x_j}$.
In other words, $\sem{\lambda( x, \alpha )}_H = \exp( iH )$ where $H = \alpha \bigotimes_{j=0}^{n-1} \ket{x_j} \bra{x_j}$.

\begin{theorem}
    \label{Thm:LambdaSem}
    If $x \in \{ 0, 1 \}^n$ and $\alpha \in \mathbb{R}$, then $\lambda( x, \alpha ) \in \Core( n,n )$ and $\sem{\lambda( x, \alpha )}_H = \exp( iH )$ where $H = \alpha \bigotimes_{j=0}^{n-1} \ket{x_j} \bra{x_j}$.
    If $\alpha \in ( -\pi, \pi ]$, then $\Log\sem{\lambda( x, \alpha )}_H = iH$ as well.
\end{theorem}

\begin{corollary}
    \label{Cor:Diag}
    Let $n \in \mathbb{N}$ with a bijection $f: \{ 1, 2, \ldots, 2^n \} \to \{ 0, 1 \}^n$.
    If $M \in \CatUnitary( 2^n, 2^n )$ is a diagonal matrix, then there exists a unique function $\alpha: \{ 1, 2, \ldots, 2^n \} \to ( -\pi, \pi ]$ such that $M = \sem{C}_H$ where $C = \lambda( f( 1 ), \alpha( 1 ) ) \circ \lambda( f( 2 ), \alpha( 2 ) ) \circ \cdots \circ \lambda( f( 2^n ), \alpha( 2^n ) )$.
\end{corollary}

\begin{proof}
    Since $M$ is a diagonal matrix, then $M$ has eigenbasis $\{ \ket{y} : y \in \{ 0, 1 \}^n \}$.
    For each $j \in \mathbb{N}$, let $\ket{b_j}$ denote $\ket{f( j )}$ and $\beta_j$ denote the eigenvalue associated to $\ket{b_j}$.
    Then define $\alpha( j ) = \Log( \beta_j ) / i$.
    It remains to be shown that $\sem{C}_H = M$.
    By \cref{Thm:LambdaSem}, $\Log\sem{\lambda( f( j ), \alpha( j ) )}_H = i \alpha( j ) \ket{b_j} \bra{b_j}$ for each $j \in \{ 1, 2, \ldots, 2^n \}$.
    Then it suffices to show that each $\Log\sem{\lambda( f(j), \alpha(j) )}_H$ commutes with each $\Log\sem{\lambda( f(k), \alpha(k) )}_H$.
    To this end, let $j \in \{ 1, 2, \ldots, 2^n \}$ and $k \in \{ 1, 2, \ldots, 2^n \}$ such that $j \ne k$.
    Since $f$ is a bijection, then $f(j) \ne f(k)$.
    Then $\braket{b_j|b_k} = \delta_{j,k} = 0 = \delta_{k,j} = \braket{b_k|b_j}$ where $\delta$ is the Kronecker delta function.
    Then the following equation holds.
    \begin{equation*}
        \Log\sem{\lambda( f( j ), \alpha( j ) )}_H\Log\sem{\lambda( f( k ), \alpha( k ) )}_H
        =
        \mathbb{O}_{2^n}
        =
        \Log\sem{\lambda( f( k ), \alpha( k ) )}_H\Log\sem{\lambda( f( j ), \alpha( j ) )}_H
    \end{equation*}
    Since $j$ and $k$ were arbitrary, then $\Log\sem{\lambda( f(j), \alpha(j) )}_H$ commutes $\Log\sem{\lambda( f(k), \alpha(k) )}_H$ for each $j \in \{ 1, 2, \ldots, 2^n \}$ and $k \in \{ 1, 2, \ldots, 2^n \}$.
    Then the following equation holds.
    \begin{align*}
        \sem{C}_H
        =
        \prod_{j=1}^{2^n} \sem{\lambda( f( j ), \alpha( j )}_H
        &=
        \prod_{j=1}^{2^n} \exp( i \alpha(j) \ket{b_j} \bra{b_j} )
        \\
        &=
        \exp\left( \sum_{j=1}^{2^n} i \alpha(j) \ket{b_j} \bra{b_j} \right)
        &
        ( \text{by~\cref{Eqn:SimulDiag}} )
        \\
        &=
        \sum_{j=1}^{2^n} \exp( i \alpha(j) ) \ket{b_j} \bra{b_j}
        \\
        &=
        \sum_{j=1}^{2^n} \exp( i \Log( \beta_j ) / i ) \ket{b_j} \bra{b_j}
        \\
        &=
        \sum_{j=1}^{2^n} \beta_j \ket{f(j)} \bra{f(j)}
        =
        M
    \end{align*}
    It remains to be shown that $\alpha$ is unique.
    Assume that there exists some $\gamma: \{ 1, 2, \ldots, 2^n \} \to (-\pi, \pi]$ such that $M = \sem{D}_H$ where $D = \lambda( f( 1 ), \gamma( 1 ) ) \circ \lambda( f( 2 ), \gamma( 2 ) ) \circ \cdots \circ \lambda( f( 2^n ), \gamma( 2^n ) )$.
    It suffices to show that $\alpha$ and $\gamma$ agree pointwise.
    Let $j \in \{ 1, 2, \ldots, 2^n \}$.
    Then the following equation holds.
    \begin{equation*}
        e^{i\alpha(j)}
        =
        \sum_{k=1}^{2^n} e^{i\alpha(j)} \delta_{j,k} \delta_{k,j}
        =
        \bra{b_j} \sem{C}_H \ket{b_j}
        =
        \bra{b_j} \sem{D}_H \ket{b_j}
        =
        \sum_{k=1}^{2^n} e^{i\gamma(j)} \delta_{j,k} \delta_{k,j}
        =
        e^{i\gamma(j)}
    \end{equation*}
    Since $\alpha(j) \in (-\pi,\pi]$ and $\beta(j) \in (-\pi,\pi]$, then the following equation holds.
    \begin{equation*}
        \alpha(j)
        =
        i\alpha(j)/i
        =
        \Log( \exp( i\alpha( j ) ) ) / i
        =
        \Log( \exp( i\gamma( j ) ) ) / i
        =
        i\gamma(j)/i
        =
        \gamma(j)
    \end{equation*}
    Since $j$ was arbitrary, then $\alpha = \gamma$.
    Since $\gamma$ was arbitrary, then $\alpha$ is unique.
\end{proof}
